
% ══════════════════════════════════════════════════════════════════════════════
%  APPENDIX A: DECLARATION OF AI ASSISTANCE
% ══════════════════════════════════════════════════════════════════════════════

\chapter{Declaration of AI Assistance}

\section{Use of Artificial Intelligence in This Document}

This course guide was developed with assistance from artificial intelligence (AI)
in the following capacities:

\begin{itemize}
  \item \textbf{Template Design:} AI provided guidance on structuring the document layout,
        choosing appropriate color schemes.

  \item \textbf{LaTeX Code Generation:} AI assisted in writing and refactoring LaTeX code,
        including package configurations, box styling (tcolorbox), TikZ decorations,
        and custom command definitions.

  \item \textbf{LaTeX Structure Refactoring:} AI helped reorganize and optimize the document
        structure for maintainability, readability, and professional presentation.

  \item \textbf{Mathematical Formula Formatting:} AI provided assistance in formatting and
        presenting mathematical equations in clear, readable notation aligned with
        conventions.
\end{itemize}

\bigskip\noindent
\textbf{Author Verification:} All content, including scientific accuracy, theoretical derivations,
example problems, and pedagogical choices, has been carefully reviewed and verified by the author.
The AI assistance was limited to technical formatting and structural optimization, not to
the generation or validation of course content itself.



\clearpage

% ══════════════════════════════════════════════════════════════════════════════
%  APPENDIX B: MATHEMATICAL TOOLS & DERIVATIONS
% ══════════════════════════════════════════════════════════════════════════════

\chapter{Mathematical Tools \& Derivations}

\section{Useful Identities}

\subsection*{Trigonometric}
$\sin^2\theta + \cos^2\theta = 1$,\quad
$\tan\theta = \sin\theta/\cos\theta$,\quad
$\sin(A\pm B) = \sin A\cos B \pm \cos A\sin B$.

\subsection*{Hyperbolic}
$\cosh^2 x - \sinh^2 x = 1$,\quad
$\sinh(x\pm y) = \sinh x\cosh y \pm \cosh x\sinh y$.

\subsection*{Common Integrals}
\begin{align*}
  \int x^n\,dx &= \tfrac{x^{n+1}}{n+1} + C \quad (n\neq-1), &
  \int e^{ax}\,dx &= \tfrac{1}{a}e^{ax}+C,\\
  \int \frac{dx}{x} &= \ln|x|+C, &
  \int \sin x\,dx &= -\cos x + C.
\end{align*}

\subsection*{Taylor Expansions}
\begin{align*}
      f(x) &= f(a) + f'(a)(x-a) + \tfrac{f''(a)}{2!}(x-a)^2 + \cdots\\
\end{align*}
\begin{align*}
      \sin x &= \sum_{n=0}^\infty \frac{(-1)^n x^{2n+1}}{(2n+1)!} = x - \tfrac{x^3}{3!} + \tfrac{x^5}{5!} - \cdots\\
      \cos x &= \sum_{n=0}^\infty \frac{(-1)^n x^{2n}}{(2n)!} = 1 - \tfrac{x^2}{2!} + \tfrac{x^4}{4!} - \cdots \\
      e^x &= \sum_{n=0}^\infty \frac{x^n}{n!} = 1 + x + \tfrac{x^2}{2!} + \cdots \\
      \sqrt{1+x} &= \sum_{n=0}^\infty \binom{1/2}{n} x^n = 1 + \tfrac{1}{2}x - \tfrac{1}{8}x^2 + \cdots
\end{align*}

\section{Differential Geometry Primer}

The \textbf{metric tensor} $g_{\mu\nu}$ encodes the geometry of spacetime:
$ds^2 = g_{\mu\nu}\,dx^\mu\,dx^\nu$.
The \textbf{Christoffel symbols} (not tensors) are defined as
\begin{equation}
  \Gamma^\lambda_{\mu\nu} = \tfrac{1}{2}g^{\lambda\sigma}
    (\partial_\mu g_{\nu\sigma} + \partial_\nu g_{\mu\sigma} - \partial_\sigma g_{\mu\nu}).
\end{equation}
The \textbf{Ricci tensor} is the contraction $R_{\mu\nu} = R^\lambda{}_{\mu\lambda\nu}$
of the Riemann curvature tensor, and the Ricci scalar is $R = g^{\mu\nu}R_{\mu\nu}$.
These enter the Einstein field equations as
$G_{\mu\nu} \equiv R_{\mu\nu} - \tfrac{1}{2}g_{\mu\nu}R = (8\pi G/c^4)T_{\mu\nu}$.

\section{Relativistic Limits}
The energy-momentum relation is $E^2 = p^2c^2 + m^2c^4$.
So in a relativistic limit, we can have:
\[
      \{ E, m, \rho\} \ll \{p, T\}
\]
which any one of the values in left-hand set is much smaller than any one of the values in the right-hand set.
\clearpage

% ══════════════════════════════════════════════════════════════════════════════
%  APPENDIX C: CONSTANTS & CONVERSION FACTORS
% ══════════════════════════════════════════════════════════════════════════════

\chapter{Constants \& Conversion Factors}

\section{Physical Constants}

\begin{center}
\begin{tabular}{lll}
\toprule
\textbf{Quantity} & \textbf{Symbol} & \textbf{Value (SI)} \\
\midrule
Speed of light         & $c$       & $2.998 \times 10^8 \,\mathrm{m\,s^{-1}}$           \\
Gravitational constant & $G$       & $6.674 \times 10^{-11} \,\mathrm{m^3\,kg^{-1}\,s^{-2}}$ \\
Hubble constant today  & $H_0$     & $67.4 \pm 0.5 \,\mathrm{km\,s^{-1}\,Mpc^{-1}}$     \\
Cosmological constant   & $\Lambda$ & $1.11 \times 10^{-52} \,\mathrm{m^{-2}}$           \\
Cosmological constant &$\Lambda$ & $1.19 \times 10^{-35} \,\mathrm{s^{-2}}$           \\
Planck constant         & $\hbar$   & $1.055 \times 10^{-34} \,\mathrm{J\,s}$                     \\
Boltzmann constant      & $k_B$     & $1.381 \times 10^{-23} \,\mathrm{J\,K^{-1}}$                    \\
\bottomrule
\end{tabular}
\end{center}

\section{Conversion Factors}

\begin{center}
\begin{tabular}{ll}
\toprule
\textbf{Conversion} & \textbf{Value} \\
\midrule
1 astronomical unit (AU) & $1.496 \times 10^{11} \,\mathrm{m}$ \\
1 light-year (ly)   & $9.461\times10^{15}\,\mathrm{m}$ \\
1 parsec (pc)       & $3.086\times10^{16}\,\mathrm{m}$ \\
1 parsec (pc)       & $3.262\,\mathrm{ly}$ \\
1 year (yr)         & $3.156\times10^{7}\,\mathrm{s}$ \\
\midrule
1 eV                & $1.602 \times 10^{-19} \,\mathrm{J}$ \\
1 electron mass ($m_e$) & $0.511 \,\mathrm{MeV}/c^2$ \\
1 electron mass ($m_e$) & $9.109 \times 10^{-31} \,\mathrm{kg}$ \\
1 proton mass ($m_p$) & $938.3 \,\mathrm{MeV}/c^2$ \\
1 proton mass ($m_p$) & $1.673 \times 10^{-27} \,\mathrm{kg}$ \\
1 neutron mass ($m_n$) & $939.6 \,\mathrm{MeV}/c^2$ \\
1 neutron mass ($m_n$) & $1.675 \times 10^{-27} \,\mathrm{kg}$ \\
\midrule
1 Solar mass ($M_\odot$) & $1.989 \times 10^{30} \,\mathrm{kg}$ \\
1 solar radius ($R_\odot$) & $6.957 \times 10^8 \,\mathrm{m}$ \\
1 Earth mass ($M_\oplus$) & $5.972 \times 10^{24} \,\mathrm{kg}$ \\
1 Earth radius ($R_\oplus$) & $6.371 \times 10^6 \,\mathrm{m}$ \\
CMB temperature ($T_\text{CMB}$) & $2.725 \,\mathrm{K}$ \\
\bottomrule
\end{tabular}
\end{center}

\clearpage
\section{SI Unit Prefixes}
\begin{center}
\begin{tabular}{lll}
\toprule
\textbf{Prefix} & \textbf{Name} & \textbf{Factor} \\
\midrule
      $10^{-9}$ & nano (n)  & $10^{-9}$ \\
      $10^{-6}$ & micro ($\mu$) & $10^{-6}$ \\
      $10^{-3}$ & milli (m)  & $10^{-3}$ \\
      $10^{3}$  & kilo (k)   & $10^{3}$ \\
      $10^{6}$  & mega (M)   & $10^{6}$ \\
      $10^{9}$  & giga (G)   & $10^{9}$ \\
      $10^{12}$ & tera (T)   & $10^{12}$ \\
\bottomrule
\end{tabular}
\clearpage
